\section{What has been formalized, and how the proof is organized}
\label{sec:formalization}

\subsection{The exact accepted statement}

The latest fixed theorem is named
\begin{quote}\ttfamily\small
OddPiPowers.L99.\allowbreak pi\_cube\_exponentUpperBound\_l99\_exact.
\end{quote}
Its conclusion is
\begin{verbatim}
ExponentUpperBound (Real.pi ^ 3) (12970477737 / 11706979)
\end{verbatim}
The predicate used by the project expands to
\begin{equation}\label{eq:leanmeaning}
\begin{split}
\forall\varepsilon>0\ \exists c>0\ \exists Q\in\NN\ 
\forall p\in\ZZ\ \forall q\in\NN,\qquad\qquad\\
q>0,\ q\ge Q\quad\Longrightarrow\quad
cq^{-(E+\varepsilon)}\le|\xi-p/q|.
\end{split}
\end{equation}
Every integer numerator and every sufficiently large positive
denominator are included; reduced fractions are not required.
For an irrational \(\xi\), this implies the usual
\(\mu(\xi)\le E\). The project also connects the predicate to
Mathlib's infinitely-often approximation predicate
\texttt{LiouvilleWith} and proves the corresponding strict-bound
corollary at \(110793/100\).

The all-fractions formulation is useful for more than convenience.
If \(\xi=a/b\) were rational, the same value would be represented
by \(p=ta\), \(q=tb\) for arbitrarily large \(t\). The right
side of \eqref{eq:leanmeaning} would then be zero while its left
side is positive. Thus this interface itself proves irrationality.
There is no loophole in which equal fractions are silently omitted.

\subsection{From an abstract transfer to the actual family}

A transfer theorem with hypotheses about small polynomials is only
one part of the task. It must be supplied with the specific
polynomials obtained from the integral construction. For the
accepted instance the dependency chain contains the following
mathematical interfaces:
\begin{enumerate}
\item The integral and local coefficient formulas define integer
polynomials, with the required support and high-order vanishing.
\item The exact determinant gives an invertible coefficient matrix
after evaluation at \(\ii\), and after division by a nonzero
common integer divisor.
\item Coefficientwise prime valuations show that the chosen integer
really divides every coefficient. The reduced Gaussian rows are
therefore integral objects, rather than formal quotients over a field.
\item Analytic bounds supply uniform height, divided-difference,
and remainder estimates for these actual rows.
\item The finite-window prime estimate gives the required
asymptotic lower bound for their actual divisor.
\item One-sided rational numerical estimates assemble an
\texttt{AuxiliaryFamily} with \(m=171,L=99\) and the displayed
rational rates. The accepted general transfer is then applied
at \(k=3\).
\end{enumerate}
The distinction between a named hypothetical family and the actual
integral family was a substantial part of the formalization work.
The final accepted statement closes that connection.

\subsection{Why a midpoint table needs a real-variable proof}

The exact combinatorial program finds cells by sorting rational
crossings of affine endpoint functions. A finite midpoint
calculation identifies a candidate order and cost in each cell.
By itself this is a test at one point. To turn it into a theorem
for all real \(z\) in the cell, the formal certificate records
floor values, endpoint rank and inverse-rank codes, and the
candidate minimum.

The proof checks adjacent endpoint inequalities and permutation
consistency. No order change is possible without an endpoint
crossing, and all crossing positions are among the cell boundaries.
The cost is checked on each gap and at the endpoints of the inner
minimization variable \(y\). This last clause prevents an open-gap
argument from overlooking a smaller cost exactly at an endpoint.
The result is a lower bound valid for every relevant \(y\) and
every real \(z\) in each open cell. The L99 table has 3004 cells
and 56 labeled endpoints before multiplicities are expanded.

This explains the difference between the very short search code
and the larger proof data. The code proposes a finite answer; the
checked certificate supplies the order and cost facts needed to
justify that answer on a continuum. The breakpoint convention
still assigns zero to the finitely many outer exceptional points,
as in Section~\ref{sec:arithmetic}.

\subsection{How finite windows prove an infinite asymptotic estimate}

The paper formula for \(\gamma\) uses a digamma difference, hence
all prime periods. The accepted fixed proof needs only a lower
rate, so it uses finitely many periods and drops the rest. For a
cell \((a,b)\), an inward shrink by a positive rational fraction
avoids endpoint ambiguities. For each retained period \(r\), the
prime number theorem gives the contribution of
\[
\frac{N_0}{r+b'}<p<\frac{N_0}{r+a'},
\qquad a<a'<b'<b.
\]
The limit of its normalized logarithmic prime product is
\(L^{-1}((r+a')^{-1}-(r+b')^{-1})\), multiplied by the
certified cell cost. Summing a fixed finite set of such estimates
preserves their eventual validity.

For L99 the relative shrink is \(\epsilon=10^{-56}\). The
symbolic loss allowance is
\[
\frac{2\cdot99^2\epsilon\cdot1000000\cdot43\cdot3004}{99}.
\]
The 32 short prefix periods are checked explicitly; the remaining
finite range up to one million periods is bounded by a rational
telescoping estimate. The compiler does not enumerate a product
over all primes or over a million periods. The prime number
theorem is used as a proved asymptotic input, and the retained
rational sum is enough to exceed the chosen lower rate.

\subsection{Why the numerical proof does not trust floating point}

The fixed height \(1801/20\) reduces the analytic comparisons to
inequalities involving logarithms and arctangents at explicit
rational arguments. A useful logarithm identity is
\[
\log x=2\sum_{r=0}^{\infty}
\frac{1}{2r+1}\left(\frac{x-1}{x+1}\right)^{2r+1},
\qquad x\ge1.
\]
Every term is nonnegative, so truncation gives a rational lower
bound. Range reduction by a known power of two keeps the series
argument small. The L99 certificates use \(4096=2^{12}\) for
the central logarithms. For the arctangent contribution, the
elementary bound \(\arctan u\ge u-u^3/3\), \(u\ge0\), is
enough at this large height.

Consequently the analytic lemma reduces to short rational
inequalities for the 28 interval types, together with proved
monotonicity and series bounds. Ordinary kernel-checked arithmetic
then proves the advertised one-sided rates. The earlier Arb
calculation remains useful for selecting parameters and checking
the size of margins, but is not an assumption of the accepted
Lean theorem.

\subsection{Acceptance evidence and its limits}

The executed acceptance report is
\begin{quote}\small\ttfamily
formalization/OddPiPowers/logs/\allowbreak
l99-acceptance-2026-10-31.json.
\end{quote}
It records 147 new dependency modules, four final signature checks,
and 15 selected transitive axiom reports, including the actual
family and final numerical theorem. The reported logical
dependencies are \texttt{propext}, \texttt{Classical.choice}, and
\texttt{Quot.sound}. No project mathematical axiom, \texttt{sorryAx},
or native-evaluation axiom occurs in the accepted audited closure.
These are local executed verification records, not a claim of
independent specialist review.

\begin{center}
\begin{tabular}{>{\raggedright\arraybackslash}p{72mm}>{\raggedright\arraybackslash}p{65mm}}
\toprule
Assertion&Status in this revision\\
\midrule
General transfer for integer \(k\ge2\)&Accepted Lean theorem\\
L99 cubic bound \(\LNNbound\)&Accepted fixed Lean application\\
Earlier P6 cubic bound \(42170551935/35026168\)&Accepted fixed Lean application\\
Twelve-layer numerical cases \(k=5,7,9\)&Manuscript argument and Arb certificates\\
L99 fourth-power comparison at scale 19&Manuscript specialization and Arb comparison\\
Uniform bound \(301\cdot8^k\)&Paper-level proof; generic Lean assembly remains\\
Classical consequence \(B(k)\)&Derived here from a cited external theorem\\
\bottomrule
\end{tabular}
\end{center}

The earlier P6 exact value is approximately
\(1203.97275360\), slightly above its optimizing Arb value
\(1203.97243921\). The same distinction between conservative formal
rates and an optimized analytic expression applies to L99. Both
earlier sources and their accepted evidence are retained.

Proof trials were serialized under the project's fixed memory
contract. Large table reductions were split into small modules
and structural append arguments; the numerical and cell bounds
were not weakened to make the checker finish. This organization
matters for reproducibility on a machine with limited memory, but
it does not change the mathematical statement or the toolchain.
This editorial revision does not modify the accepted Lean sources
or claim new compiler acceptance for the additional exposition.
